\title{\textbf{Detecting short gravitational-wave bursts in real LIGO noise}\\[0.6em]
\large\itshape Frequentist, Bayesian and machine-learning analysis of simulated bursts injected into GWOSC O1 data}
\author{Giacomo Moretti\\[0.2em]\small Statistics and Data Analysis, Università di Milano-Bicocca, A.Y.\ 2025--2026}
\date{October 2026}
\maketitle

\begin{abstract}
\noindent We use public LIGO strain from the first observing run (O1) as a source of \emph{real}, non-Gaussian, non-stationary noise, inject simulated bursts at a known optimal SNR, and compare three ways of detecting them at the same false-alarm probability: a frequentist matched-filter bank (profile-likelihood ratio), a Bayesian analysis of the same likelihood (exact grid posterior, MCMC, Bayes factors) and machine-learning detectors (feature-based classifiers, convolutional networks, an autoencoder, a Transformer and a GRU). All methods share the data, the injections and a split by chunk of data.
(i) The whitened noise is Gaussian in the bulk but has a glitch floor of a few $10^{-3}$ per second and large epoch-to-epoch variations; they dominate the cost of low false-alarm probabilities and make the calibration of a threshold epoch-dependent.
(ii) With three phenomenological burst morphologies, the matched-filter bank detects sine-Gaussians with 50\,\% efficiency at an optimal SNR of $\approx5$ at a false-alarm probability of $10^{-2}$ per one-second window (an oracle that knows the parameters: $3.0$), and the Bayes factor detects the same bursts, with calibrated posteriors and an explicit prior dependence.
(iii) A small CNN reaches parity with the bank and is better on the shapes it was trained on; the largest effect on it is \emph{fresh injections at every epoch}, which make a single network better than a five-member ensemble. A Transformer and a GRU detect worse than the CNN, and a tuning study --- one lever at a time --- shows why: the patch embedding, not the data or the optimisation, limits the Transformer, and a convolutional stem recovers most of the gap but not all. Morphology classification is as good with fresh-injection CNNs as with Bayesian model comparison ($0.98$ against $0.97$ three-way accuracy at $\snr=8$); a goodness-of-fit of the matched filter and the energy score of the network flag signals that belong to none of the families.
(iv) With realistic binary-black-hole (IMRPhenomD) signals the banks built on bursts lose $2$--$3$ SNR units, a network trained with such signals beats an IMR matched-filter bank on waveforms never seen in training ($4.45$ against $5.40$), and with two detectors, coincidence and time slides push the false-alarm floor from $10^{-3}$ to $10^{-6}$ per window at a network SNR of $\approx11$ for balanced injections and $\approx15$ with the real antenna patterns.
(v) A second round of studies puts error bars on these statements: the seed scatter of the CNN ($0.4$ in $\snr_{50}$) is much larger than its test-set error ($0.1$), the CNN advantage lives in the tails of real noise and inside its training prior (a spectral line or an out-of-prior morphology breaks it), the public data-quality flags recover most of the gap for the matched filter, and the same pipeline transfers to O2 and O3 without retraining. Parameters of binaries are recovered to $\approx7\,\%$ in chirp mass (calibrated intervals only with exact-model MCMC), the sensitivity is expressed in $h_{\rm rss}$ and in Euclidean range ($1.3$--$1.6$\,Gpc), and the limits of one-second windows (low masses, neutron stars, NSBH, eccentricity) are quantified.
\end{abstract}
\clearpage
